%!TEX program = lualatex
\documentclass[ngerman,headheight=70pt]{scrartcl}
\usepackage[ngerman]{babel}
\usepackage{blindtext}
\usepackage[utf8]{luainputenc}
\usepackage{scrlayer-scrpage}
\usepackage[T1]{fontenc}
\usepackage{fontspec}
\usepackage{multicol}
\usepackage{enumitem}
\usepackage[german=quotes]{csquotes}
\usepackage[bottom=2cm,footskip=8mm]{geometry}
\setmainfont{Source Sans Pro}
\MakeOuterQuote{"}

\setlength{\parindent}{0pt}
\setlength{\parskip}{6pt}

\pagestyle{plain.scrheadings}
\thispagestyle{scrheadings}
\lohead{\vspace*{-1.5cm}\includegraphics[scale=0.4]{../images/stupa_logo_rot.pdf}}
\rohead{\\Präsidium des\\Studierendenparlaments}

\lofoot[Universität Hamburg · Präsidium des Studierendenparlaments\\\hrule Von-Melle-Park 5 · D-20146 Hamburg]{Universität Hamburg · Präsidium des Studierendenparlaments\\\hrule Von-Melle-Park 5 · D-20146 Hamburg}
\cofoot[]{}
\cohead[\thepage]{}

\renewcommand{\thesubsection}{TOP \arabic{subsection}}
\renewcommand{\thesubsubsection}{\arabic{subsubsection}.}
\setcounter{subsection}{-1}

\begin{document}
    UHH · StuPa-Präsidium · Von-Melle-Park 5 · D-20146 Hamburg

    %%% PAGE ONE (EDIT LATER) %%%%
    % Hier den Titel des Protokolls eintragen
    \section*{Protokoll der 3. Sitzung des Studierendenparlaments vom 26. Mai 2016}

    % Protokollanten eintragen
    \textbf{Protokoll: Jim Martens}\\
    \textbf{Ort: S08 in der ehemaligen HWP}\\
    \textbf{Beginn: 18.39 Uhr}\\
    \textbf{Ende: Uhr}

    % Wohl die komplizierteste Stelle im Protokoll
    % Als Faustregel sollten nicht mehr als 3 Namen auf einer Zeile stehen
    % Einfach an dieser Vorlage orientieren
    % Im Idealfall ist hier nicht viel zu ändern, da kaum jemand fehlt.
    % Rücktritte müssen natürlich pro Protokoll angepasst werden
    \vspace{0.5cm}
    \begin{tabular}{ll}
        Anwesend: & \\
            RCDS (5 Sitze): & Ramon Weilinger, Antonia Niecke, Ramin Shakiba, \\
                            & Jennifer Maack \\
             CampusGrün (14 Sitze): & Laura Franzen, Geoffrey Youett, Philipp Droll, \\
                                   & Tahnee Herzig, Svenja Horn, Jim Martens,\\
                                   & Mirzo Ulugbek Khatamov, Armin Günther, Martin Sievert\\
                                   & Elena Rysikova, Yasemin Günther \\
             Bier-Liste (2 Sitze) : & Jakob Pape, Karen Martirosian \\
             WiWi (2 Sitze): & Claas-Friso Hente \\
             Unicorns (5 Sitze): & Johannes Peplow, Annkathrin Löffler, \\
                                 & Andreas Hartkamp, Marielle Hermstrüwer \\
             Liste LINKS (3 Sitze): & Gunhild Berdal, Till Petersen, Sinah Mielich \\
             HWP (2 Sitze): & Ajdina Karahasan \\
             MIN (4 Sitze): & \\
             SDS* (3 Sitze): & Mena Winkler, Jacob Petersein, Artur Brückmann \\
             Bart-LISTE (2 Sitze): & Timo Zeimet, Dominic Laumer \\
             LHG (1 Sitz): & \\
             harte zeiten (1 Sitz): & Tobias Berking \\
             Jura (1 Sitz): & \\
             AL (2 Sitze): & Karima Schulze, Henri Weber \\
            & \\
        Entschuldigt: & Benjamin Welling (RCDS),\\
                      & Melf Johannsen, Mario Moldenhauer (CG), Samet Gunay,\\
                      & Kay Zöllmer (WiWi), Katharina Kucza (Unicorns), Ailina Salten, \\
                      & Lotte Rullkötter, Thea Wahlers (alle MIN), Tobias Heisig (LHG)\\
                      &\\
        Unentschuldigt abwesend: & Jan Detampel (MIN), \\
                                & Johann Baumhoefener (Jura)\\
                                &\\
        Rücktritte: & Freya Schmitz (CG) \rightarrow Sandra Fietz\\
    \end{tabular}
    \newpage
    %%% Vorgeschlagene Tagesordnung %%%
    \underline{Vorgeschlagene Tagesordnung}
    \begin{enumerate}[label={\textbf{Top \theenumi}},leftmargin=*]
        % folgende Zeile bei Sitzungsfortsetzungen einkommentieren und x durch
        % die Zahl des zuletzt behandelten TOPs ersetzen
        %\setcounter{enumi}{x}
        \item Keine Werbung von/für Bundeswehr (V1617-023) (60 Min.)
        \item Akkreditierung (V1617-024) (60 Min.)
        \item Verlegung der StuPa-Wahl in das Sommersemester (V1617-015) (60 Min.)
        \item Prüfauftrag: Wahl-O-Mat bei StuPa-Wahlen (V1617-022) (30 Min.)
        \item Verschiedenes
    \end{enumerate}

    \newpage

    % Hier beginnt Teil A. Bei Fortsetzungen kommt hier nur der Teil über
    % die Feststellung der endgültigen Fassung der Tagesordnung.
    % Für sonstige Sitzungen ist dies bereits das Gerüst für die Formalia.
    % Man muss nur noch die protokollierten Inhalte eintragen.

    {\Large \underline{Teil A}}

    \subsection{Formalia}

    Gunhild eröffnet Sitzung um 18:39.

    \subsubsection{Geschäftsbericht Präsidium}

    Freya Schmitz ist zurückgetreten.

    Katharina Fegebank kommt wahrscheinlich am 30. Juni um 19 Uhr.

    StuPa-Sitzungen am 16. und 30. Juni.

    Gunhild berichtet über Präsidiumsstand zu Ersetzungsanträgen.

    \subsubsection{Anfragen an das Präsidium}

    Droll sagt, dass es am 9. Juni ein Freiraumfestival gibt. Das ist
    problematisch.

    Till merkt an, dass man die StuPa-Sitzung in den Festivaltag
    integrieren könnte.

    Till sagt, dass man Protokolle kürzer fassen könnte und auf die Videos
    verweisen könnte.

    Keine weiteren Fragen.

    \subsubsection{Geschäftsbericht AStA}

    Droll, Franzi und Timo halten Geschäftsbericht.

    Es gab Gespräche mit Kanzler und weiteren Leuten. Dabei war mehr
    Anerkennung als Anecken vorhanden. Dabei kam es schnell zur Verständigung,
    sodass die Freiräume das Festival durchführen können.
    Auch Haushaltsplan des Café Knallhart wurde besprochen.

    AStA 0029 könnte zum Gremienraum werden. Regelung für FSRe im Philturm
    könnte auf andere Gebäude übertragen werden. WiWi-Bunker soll wohl
    schon sicher sein. Auch VMP 8 (Erzwiss) angedacht.

    Nutzungsberechtigung für Philturm nur bis Ende 2016 ausgegeben.
    Baubeginn für Philturm könnte nicht Anfang 2017 sein. Es wird wohl
    Möglichkeiten geben Räume in den alten Postgebäuden anzumieten.

    In 1,5 Wochen gibt es Gespräche mit dem Architekten der Uni.
    Vorplatz des ESA soll umgebaut werden. Schuldenuhr soll weiterhin
    stehenbleiben. Diskussion über Ästhetik mit Kanzler.

    Kampagne gegen ExIni geht weiter und viele Unterschriftenlisten aus
    den Fachbereichen zurückgekommen. Einigung bis zum 16.6. zu ExIni
    unwahrscheinlich. Nächste Woche am 2.6. VV über ExIni in Phil 6.

    In der letzten VL-Woche soll es eine Konferenz zu HealthEconomics geben.
    Die Veranstaltungen im ESA-Flügel sollen ausfallen. Ermäßigter Preis
    für Studierende beträgt 300€.

    Eine Veranstaltung mit Frau Leopold kommt nicht zustande, weil sie
    keine Zeit hat.

    Es gibt eine BA/MA-Reform in der KMK. Das erste Jahr soll notenfrei
    sein.

    Timo berichtet über die Finanzen. Nächste Woche soll der erste Entwurf
    AStA-intern vorgelegt werden.

    Geoffrey hält Bericht aus dem Referat ÖffKulturHoPo. Es wurde Postverteiler
    aufgebaut, über den den FSRen alle 2 Wochen Material zugeschickt werden soll.
    Es wird sich Gedanken gemacht, wie man die FSRe besser in die inhaltliche
    Arbeit und Demos einbeziehen kann.

    Es soll Fotoaktion zur Frage geben "Was bedeutet Uni für dich?". Es wurde
    sich überlegt, was an Veranstaltungen zur EM gemacht werden kann.

    Es gab immer Frage, wie Campus OpenAir gestaltet werden soll. Freiraum-Festival
    ist im Prinzip genau das: Ein studentisch organisiertes Festival ohne externe
    Firma. AStA als Kooperationspartner auf Plakaten drauf. Ab 16 Uhr beginnt
    das Festival. Ab 18 Uhr wird mehr Musik gemacht bis 22 Uhr. Evtl. Sondergenehmigung
    um bis 24 Uhr Musik machen zu können. Es wird auch Stände geben

    AStA hat sich überlegt, wie er sich dort einbringen kann. Idee ist sich mit
    EM-Veranstaltungen einzubringen.

    19:14

    Kulturkurse laufen alle. Es gibt einige Kinderkrankheiten. Es soll das Buchungssystem
    optimiert werden und eine neue Website gemacht werden. Im nächsten Semester soll
    es noch besser und fruchtbarer ablaufen.

    Jacob berichtet aus Antidis/Anfifa/Krit. Wiss. Referat kann bislang
    nicht wirklich durchstarten, da viele Personen nicht da sind.

    Übergreifendes Programm sind EM-Veranstaltungen. Das Ziel ist es die
    Länder der gegnerischen Mannschaften vorzustellen und parallel zu
    den Spielen ein OpenAir-Kino anzubieten. Überlegung zu verschiedenen
    Veranstaltungsarten. Am 9.7. soll es ein Fußballturnier in Kooperation
    mit Refugees Welcome Karoviertel stattfinden.

    Die Veranstaltung zum NS-Studentenbund war gut besucht und hat eine
    heterogene Zuschauermenge angezogen. Jacob fasst die Veranstaltung
    zusammen.

    Artur berichtet vom AS. Vertreter für den HR wurden gewählt. Nöth
    wurde in den HR gewählt. Platzer macht weiter. Bei Ziel- und Leistungsvereinbarungen
    hat sich Unileitung und Behörde darauf geeinigt die quantitativen Kriterien
    niedrig anzusetzen, sodass diese tatsächlich erfüllbar sind. Dies sei der
    erste Schritt auf dem Weg zu qualitativen Kriterien.

    AS beschließt mit 16:0:2, dass der Aufruf der Wissenschaftler*Innen
    für Frieden unterzeichnet wird.

    Mit dem Studierendenwerk wurde vereinbart, dass die WissenschaftsA-Sitzung
    gemeinsam vorbereitet wird.

    Am 30.5. findet im ESA B die letzte Veranstaltung der Reihe "ich konsumiere,
    du konsumierst, wir kritisieren" statt.

    Es wird eine Veranstaltung zum Thema ExIni geben - zusammen mit WiMis
    der Uni Jena.


    Olli berichtet aus dem Ausschuss gegen Rechts. Am 5.7. wird es den Tag
    gegen Rechts geben.

    \subsubsection{Anfragen an den AStA}

    Till interessiert sich dafür, was für ein Termin die Ausschusssitzung
    mit dem Studierendenwerk ist.

    Franzi erklärt, wie es zur Sitzung gekommen ist. Artur ergänzt.

    19:30

    Vor dem Hintergrund "Form follows Function" fragt Till, inwieweit
    das ÖffRef die Form bzgl. des zu verbreitenden Inhalts reflektiert
    hat.

    Till erwähnt, dass nicht alle FSRe mit der Regelung im Philturm
    zufrieden sind. Von daher ist die Verallgemeinerung dessen nicht
    per se sinnvoll. Zum Beispiel die FSRE im Erzwiss müssen eingebunden
    werden.

    An Geoffrey geht die Frage, was mit dem Begriff Freiraumkultur gemeint
    ist.

    Geoffrey beantwortet die Frage von Till zu den Formen. Der Fokus
    bei Veranstaltungen soll der persönliche Kontakt sein.

    Das Referat ist erst spät in die Diskussion um das Festival eingestiegen.
    Es gab Einigkeit bei Gesprächen, dass es politischen Hintergrund gibt, dass
    Freiräume existieren. Man hat sich in Bezug zum Festival vorgenommen,
    diese Diskussion was Freiraumkultur ist in den folgenden Monaten
    fortzuführen.

    Franzi findet die Regelung im Philturm auch nicht ideal. Sie sei ein
    Zwischenstand. Zwischen den FSRen verläuft die Konfliktlinie, wofür längerer
    Zugang benötigt wird. Die einen wollen arbeiten, die anderen Partys feiern.
    In Bezug zu VMP 8 gibt es Gespräche.

    19:47

    Feststellung Beschlussfähigkeit: 25/47

    \subsubsection{Dringlichkeitsanträge des AStA}

    Keine.

    \subsubsection{Aktuelle Stunde (falls entsprechender Antrag vorliegt)}

    Keine.

    \subsubsection{Feststellung der endgültigen Fassung des Teils B der Tagesordnung}

    Jacob möchte TOP 4 und 3 tauschen. 15:3:4

    Tagesordnung wurde angenommen ohne Widerspruch.

    %%% BESCHLOSSENE TAGESORDNUNG %%%
    \underline{Endgültige Fassung der Tagesordnung}
    \begin{enumerate}[label={\textbf{Top \theenumi}},leftmargin=*]
        % folgende Zeile bei Sitzungsfortsetzungen einkommentieren und x durch
        % die Zahl des zuletzt behandelten TOPs ersetzen
        %\setcounter{enumi}{x}
        \item Keine Werbung von/für Bundeswehr (V1617-023) (60 Min.)
        \item Akkreditierung (V1617-024) (60 Min.)
        \item Prüfauftrag: Wahl-O-Mat bei StuPa-Wahlen (V1617-022) (30 Min.)
        \item Verlegung der StuPa-Wahl in das Sommersemester (V1617-015) (60 Min.)
        \item Verschiedenes
    \end{enumerate}

    \subsubsection{Feststellung der Beschlussfähigkeit}

    Schon geschehen.

    \subsubsection{Genehmigung der Protokolle der vorangegangenen Sitzungen}

    Protokoll vom 1. Fortsetzung mit einer Enthaltung und 0 Gegenstimmen
    angenommen.

    Protokoll 2. Sitzung bei 3 Enthaltungen angenommen.

    20:00

    Pause bis 20:10, verlängert bis 20:15

    \vspace{0.5cm}
    {\Large \underline{Teil B}}

    % Teil B beginnt
    % Dieser wird sich für jede Sitzung unterscheiden.
    % Jeder TOP ist eine subsection.
    % Vorhandene Beispiele anfangs entfernen und durch tatsächliche TOPs ersetzen

    % Folgende Zeile bei Sitzungsfortsetzungen einkommentieren und x mit
    % der Zahl des zuletzt behandelten TOPs ersetzen.
    %\setcounter{subsection}{x}
    \subsection{Keine Werbung von/für Bundeswehr}

    Franzi stellt den Antrag vor.

    Tobias möchte, dass die beiden Absätze des Antrags getrennt werden
    und eigenständig behandelt werden. Außerdem hat er einen Änderungsantrag:

    Zeile 4 von Universität bis Zeile 5 werden ersetzen durch:
    Wissenschaft hat die Aufgabe, die Wahrheit für die menschenwürdige
    Entwicklung der Welt zu bestimmen und zur Wirklichkeit zu machen. Die
    Wahrheit stirbt jedoch im Krieg bekanntlich zuerst.

    Artur hält einen Redebeitrag und erklärt den historischen Zusammenhang
    der Gründung der Uni Marketing GmbH.

    20:32

    Gunhild möchte "künftigen" in Zeile 7 streichen. Am 11. Juni findet
    der Tag der Bundeswehr statt. Gunhild befürwortet die Teilung des Antrags.

    Ramin stellt GO-Antrag auf Feststellg. der Beschlussfähigkeit. Till beantragt
    eine Fraktionspause für Liste Links von 20 Minuten bis 21:03.

    In dieser Zeit wird weiter diskutiert.

    Geoffrey hat eine persönliche Erklärung abgegeben:

    Sinnbildlich spiegeln die gähnend leeren Sitzplätze der Opposition
    wider, was hinter ihren Wahlprogrammen steckt -- nichts.

    Jetzt folgt die Abstimmung: Mit 22 Parlamentarier*Innen nicht beschlussfähig.

    Sitzung ist abgebrochen um 21:08. Nächste Sitzung am 9. Juni.
\end{document}
